\centerline{\bf \Huge Октябрьский зачёт по алгебре }

\begin{flushright}
        \scriptsize{\textbf{Разве ты не заметил, что способный к математике изощрен во всех науках в природе?\\
         Платон}}
        
\end{flushright}

\problem{\textit{\textbf{1}}} Вывод формулы корней квадратного уравнения через дискриминант.

\problem{\textit{\textbf{2}}} Вывод Теормы Виета. (Если доказываете, используя формулу корней через дискриминант, но саму формулу не доказали => такой вывод не засчитывается)

\begin{flushright} \textbf{3 балла}\\ 
\end{flushright}

\problem{\textit{\textbf{3}}} гЕниальная задача: $1^2=1, 2^2=4,$ продолжите до $,100^2=10000$.

\problem{\textit{\textbf{4}}} Вывод частного случая Теоремы Безу. (Если доказываете, используя формулу корней через дискриминант, но саму формулу не доказали => такой вывод не засчитывается)

\begin{flushright} \textbf{6 баллов}\\ 
\end{flushright}

\problem{\textit{\textbf{5}}} Вычислите, раскрыв скобки: $(a-b)^{17} - (a+b)^{17}=?$

\problem{\textit{\textbf{6}}} Докажите, что $\sqrt{2025}$ - рациональное число, а $\sqrt{2026}$ - иррациональное число.

\problem{\textit{\textbf{7}}} Докажите, что множество чётных, целых, рациональных чисел - это счётные множества, а множество иррациональных чисел - несчётно. (Биекция? Мощность множества? Равномощные множества? Счётное множество?)

\begin{flushright} \textbf{9 баллов}\\ 
\end{flushright}

\centerline{\textit{Упростите сложные радикалы:}}

\problem{\textbf{1}}  $\sqrt{|12\sqrt{5} -  29|} - \sqrt{12\sqrt{5} +  29}=?$

\problem{\textbf{1!}} $\sqrt{13-4\sqrt{14-2\sqrt{22-\sqrt{321+\sqrt{1280}}}}}=?$

\hspace{3mm}

\centerline{\textit{Избавьтесь от иррациональности в знаменателе и преобразуйте по возможности:}}

\problem{\textbf{2}} $\dfrac{\sqrt{\sqrt{3697} + \sqrt{993}}}{\sqrt{\sqrt{3697} - \sqrt{993}}}$

\problem{\textbf{2!}} $\dfrac{25}{\sqrt{3 + \sqrt{2}}-\sqrt{2}} \cdot \dfrac{46}{3+\sqrt{2}-\sqrt{1 + \sqrt{2}}}$

 \newpage


\centerline{\textit{ Сравните числа:}}

\problem{\textbf{3}} $\sqrt{37}$ vs $\sqrt{77}-\sqrt{7}$

\problem{\textbf{3!}} $-2024+\sqrt{\sqrt{17+12\sqrt{2}}}$ vs $\sqrt{2}+\sqrt{3}-2025$

 \hspace{3mm}

\centerline{\textit{ Решите уравнения:}}

\problem{\textbf{4}} $9x^2-24x-|3x-4|=4$

\problem{\textbf{4!}}  $\left(\dfrac{5|x|+1}{2|x|-3}\right)^2 + \left(\dfrac{3-2|x|}{5|x|+1}\right)^2 = \dfrac{82}{9}$

 \hspace{3mm}


\problem{\textbf{5}} $\dfrac{30}{x^2-1} + \dfrac{7-18x}{x^3+1}=\dfrac{13}{x^2-x+1}$

\problem{\textbf{5!}} $\dfrac{4x^2}{8x^6+1} + \dfrac{1}{16x^8-4x^4+4x^2-1}=\dfrac{2}{4x^4+2x^2-1}$

 \hspace{3mm}

 \problem{\textbf{6}} Найдите произведение наименьшего и наибольшего целого значения параметра $e$, при котором уравнение:
 
 $25x^2-10ex+2025=0$ не имеет действительных корней.
 
 \problem{\textbf{6!}} При каких значениях параметра $t$ уравнения равносильны? 

$x^2+2(t-3)x+(t^2-7t+12)=0$ 

$x^2+(t^2-5t+6)x=0$

\textit{Равносильны => Множества решений обоих уравнений полностью совпадают.}

 \problem{\textbf{7}} В квадратном уравнении сумма третьих степеней корней равна 10. Найдите $l$.

 $25x^2-10x-201l^2=0$

 \problem{\textbf{7!}} Известно, что корни уравнения $x^2-13x+i=0$ равны соответственно квадратам корней уравнения $x^2+tx+6=0$. Найдите $i$ и $t$ и корни каждого из уравнений.

  \hspace{3mm}

\centerline{\textbf{1,2,3 - 6} баллов; \textbf{4,5 - 9} баллов; \textbf{6,7 - 12} баллов}

\centerline{\textbf{1!,2!,3! - 12} баллов; \textbf{4!,5! - 18} баллов; \textbf{6!,7! - 24} балла}



 